
O UV mapping é uma técnica essencial para associar informações de textura a superfícies tridimensionais. A partir do mapeamento tradicional, transformações aplicadas às coordenadas UV permitem criar comportamentos dinâmicos, enquanto métodos procedurais possibilitam calcular padrões e resultados por meio de funções matemáticas.

A revisão bibliográfica mostrou que o texture mapping possui fundamentos consolidados, incluindo problemas geométricos e de filtragem, enquanto a geração procedural amplia as possibilidades de criação e reutilização de conteúdo. As publicações analisadas também mostram que métodos procedurais podem ser úteis para produzir grandes quantidades de conteúdo, mas exigem atenção ao controle artístico e ao custo computacional.

No contexto de Vulkan, a técnica pode ser integrada ao pipeline gráfico por meio de buffers de vértices, shaders, imagens, image views e samplers. As coordenadas UV podem ser interpoladas durante a rasterização e transformadas antes da amostragem da textura, permitindo efeitos dinâmicos e procedurais.

As aplicações práticas incluem superfícies animadas, materiais procedurais, terrenos, ambientes virtuais e efeitos visuais. Os casos de Rocket League e Marvel's Spider-Man Remastered foram escolhidos para relacionar os conceitos teóricos ao resultado visual observado em jogos.

Para a versão final da entrega, as capturas de gameplay deverão ser substituídas pelos registros próprios do grupo e os vídeos deverão ser adicionados somente quando uma captura estática não for suficiente para demonstrar o comportamento analisado.
